\section{Methodology}
\label{sec:methodology}

\noindent\textbf{Note on review classification.}
This review was conducted by a single author and accordingly cannot satisfy the dual-screener requirement of PRISMA 2020.
It is classified throughout as a \emph{structured literature review} rather than a PRISMA-compliant systematic review.
PRISMA-format elements (flow diagram, exclusion counts, search strategy reporting) are retained for transparency and reproducibility, with the caveat that the corpus was assembled under single-author constraints.
The implications of this for the review's headline finding (five of 237 claims verified) are quantified in the sensitivity analysis in Section~\ref{subsec:verification_protocol}.

We follow a structured literature review protocol drawing on the guidelines of
Kitchenham and Charters~\cite{kitchenham2007guidelines} and the PRISMA 2020
statement~\cite{page2021prisma}, adapted for single-author execution as described below.

\subsection{Research Questions}

The review is structured around four primary research questions (RQs):

\begin{description}
  \item[RQ1] What architectural approaches are used to implement guardrails for LLMs,
             and how do they differ in placement, implementation layer, and modality
             coverage?
  \item[RQ2] What mechanisms have been proposed for real-time, inference-time, and
             streaming guardrails, and what are their documented latency and accuracy
             characteristics?
  \item[RQ3] What safety challenges are specific to multimodal LLMs, and what
             defences have been proposed?
  \item[RQ4] What are the principal reproducibility limitations and open problems
             in the current guardrail literature?
\end{description}

\subsection{Search Strategy}

We searched the following databases: arXiv (cs.AI, cs.CL, cs.CR, cs.LG), ACL
Anthology, IEEE Xplore, ACM Digital Library, USENIX Security proceedings, and
the NeurIPS/ICML/ICLR proceedings (2022--2025).
The review window is \textbf{January 2022 to 30 June 2025} (inclusive); all
corpus papers have an arXiv submission date or formal publication date that falls
within this window.
This cutoff is applied consistently throughout the paper; the single entry
submitted in January 2025 (Aegis2.0, arXiv:2501.09004) falls within the window.
The lower bound of 2022 is motivated by the emergence of instruction-tuned LLMs
as production systems: the deployment of InstructGPT~\cite{ouyang2022training}
and the public release of ChatGPT in late 2022 mark the point at which guardrails
became a practical engineering requirement rather than a theoretical concern.
Papers published before 2022 that are cited in this review (Weidinger~et~al.\
2021~\cite{weidinger2021ethical}; Bommasani~et~al.\ arXiv v1 August
2021~\cite{bommasani2022opportunities}; Perez~et~al.\ 2022~\cite{perez2022red};
FUDGE 2021~\cite{yang2021fudge}; NeuroLogic 2021~\cite{lu2021neurologic}) are
background or adjacent-field references, not primary corpus members; the corpus
composition note in Section~\ref{sec:methodology} (Data Extraction subsection)
classifies them accordingly.
Search strings were constructed from three concept clusters and executed as follows.

\paragraph{Primary Boolean query string}
Applied across ACM DL, IEEE Xplore, Springer Link, and arXiv:
\begin{quote}\footnotesize\sloppy
(``large language model'' OR ``LLM'' OR ``foundation model'')
AND (``guardrail'' OR ``safety filter'' OR ``content moderation'' OR ``output filter'' OR ``alignment'' OR ``safety mechanism'')
AND (``real-time'' OR ``streaming'' OR ``latency'' OR ``token-by-token'' OR ``multimodal'' OR ``vision-language'' OR ``audio'' OR ``video'' OR ``jailbreak'' OR ``adversarial'')
\end{quote}

\paragraph{Secondary query (broadening for coverage)}
\begin{quote}\footnotesize\sloppy
(``LLM safety'' OR ``AI safety'' OR ``content filtering'')
AND (``runtime'' OR ``inference-time'' OR ``deployment'')
AND (``2022'' OR ``2023'' OR ``2024'' OR ``2025'')
\end{quote}

\paragraph{arXiv-specific mechanism}
arXiv was searched via the web interface (\texttt{arxiv.org/search}) using the cs.AI, cs.CL, and cs.CR category filters with the primary query applied to title and abstract fields.
The arXiv API was not used; searches were conducted manually with results sorted by relevance and filtered by submission date (January 2022 -- June 2025).
Readers attempting to replicate this search should note that the arXiv web interface does not produce a stable, versioned query URL; result ordering and coverage may differ across sessions due to index updates, and exact replication of the retrieved set is not guaranteed from the interface alone.

\paragraph{Search execution dates}
All database searches were conducted in May--June 2025.
ACM DL: June~1, 2025.
IEEE Xplore: June~2, 2025.
Springer Link: June~3, 2025.
arXiv: June~4--7, 2025 (multiple passes across cs.AI, cs.CL, cs.CR).
USENIX proceedings: June~8, 2025.

\paragraph{USENIX inclusion justification}
USENIX Security and USENIX SOUPS were included because they represent the primary publication venues for security-oriented LLM research that does not appear in IEEE or ACM databases; prior surveys (Dong et al.\ 2024, Yao et al.\ 2023) similarly include USENIX venues.

\paragraph{Snowball stopping criterion}
Forward and backward citation snowball was applied to all included papers.
Snowball was terminated when a complete forward/backward pass yielded no citation not already present in the corpus.
This required three rounds; the third round added zero new papers.

The three concept clusters underlying the primary query were:

\begin{enumerate}
  \item \textit{Target}: ``large language model'' OR ``LLM'' OR ``foundation model''
        OR ``vision-language model'' OR ``VLM'' OR ``multimodal''
  \item \textit{Mechanism}: ``guardrail'' OR ``safety'' OR ``alignment'' OR
        ``moderation'' OR ``content filter'' OR ``jailbreak''
  \item \textit{Constraint}: ``real-time'' OR ``inference-time'' OR ``streaming'' OR
        ``latency'' OR ``online'' OR ``multimodal'' OR ``cross-modal''
\end{enumerate}

Strings were combined with AND across clusters.
One term in the Constraint cluster warrants a note on search precision: the term
``online'' is a high-recall, low-precision choice that matches papers on online
learning, online platforms, and real-time systems more broadly, many of which are
off-topic.
Its inclusion was motivated by the desire to capture all papers framing safety as
an online decision problem; the precision cost is visible in the exclusion counts
below (510 records excluded for ``no guardrail mechanism'' at the title/abstract
stage, a count inflated by off-topic ``online'' matches).
Estimated title/abstract precision is approximately 19\% (355 advanced to full
text from 1,847 screened), which is within the typical range for broad SLR searches
in fast-moving fields but confirms that the search string was tuned for recall
rather than precision.

Database search returned 1,812 records.
Forward and backward reference snowballing was applied to all included papers,
adding a further 35 records (total identification: 1,847; see Figure~\ref{fig:prisma}).
\begin{sloppypar}
The PRISMA 2020 exclusion-stage breakdown is as follows (figures also encoded
in the flow diagram).
\textit{Title/abstract screening}: 1,847 records screened;
1,492 excluded (pre-LLM NLP safety $n=412$; no guardrail mechanism $n=510$;
pure policy/legal $n=298$; duplicate/superseded $n=272$);
355 advanced to full-text review.
\textit{Full-text screening}: 355 assessed;
280 excluded (no empirical/formal evaluation $n=134$;
not novel $n=89$; inaccessible $n=31$; out-of-scope $n=26$);
75 studies included in the final corpus.
\end{sloppypar}

\subsection{Inclusion and Exclusion Criteria}

\paragraph{Included} Peer-reviewed papers and arXiv preprints reporting: (i) novel
guardrail mechanisms for LLMs; (ii) attacks against LLM safety mechanisms; (iii)
evaluation frameworks or benchmarks for LLM safety; (iv) surveys or analyses of
LLM safety with primary empirical content.

\paragraph{Excluded} Papers focused exclusively on classical NLP safety (pre-LLM
era); pure policy or legal analyses without technical content; papers without any
empirical or formal evaluation; duplicates (earlier version superseded by a later
published version).

\subsection{Screening Process}

Title and abstract screening was performed by the single author.
Full-text screening was applied to all papers passing abstract screening.
PRISMA 2020 requires two independent screeners to establish inter-rater reliability
(IRR); this requirement is not met and cannot be satisfied post-hoc by any
intra-author measure.
\textbf{This is a structural defect in the methodology, not a cosmetic limitation.}
The absence of a second independent screener means that systematic biases in
inclusion judgements — for instance, a tendency to include papers whose abstracts
pattern-match to the author's prior reading — cannot be detected or corrected.
No intra-author procedure can substitute for genuine inter-rater agreement.
This review is accordingly classified as a structured literature review rather than a PRISMA-compliant systematic review;
all PRISMA-format elements (flow diagram, exclusion counts) are retained for transparency,
with the understanding that the corpus was assembled under this constraint.

To mitigate the structural disadvantage of single-author screening, the following procedural compensations were applied:
(i) a pre-specified codebook with operationalised inclusion/exclusion criteria was finalised before screening began and was not modified during screening;
(ii) an audit trail was maintained in a versioned extraction spreadsheet recording the inclusion decision and primary reason for each screened paper;
(iii) borderline inclusion decisions were resolved by consulting the primary literature of the citing paper (i.e., if a borderline paper was cited by a clearly-in-scope paper for a specific claim, the borderline paper was included);
(iv) the claim verification protocol used three temporally separated passes with mandatory unanimous agreement, which functions as a within-author reliability check on the claim extraction.

To document the degree of intra-author consistency (not as a substitute for IRR,
but as a transparency disclosure), test-retest reliability was assessed on a
10\% random sample ($n = 22$ papers) in which the author conducted a blinded
re-screening four weeks after the initial pass.
Agreement on inclusion/exclusion was 91\% (20/22 papers; 95\% Wilson score
confidence interval 71\%--99\%); the two discrepant cases were resolved by
consulting a third-party domain expert.
This figure is reported as \emph{test-retest reliability only}: because all
passes were conducted by the same author, Cohen's $\kappa$ — which formally
requires two distinct raters — is undefined and is not reported.
The 91\% raw agreement rate \emph{does not} constitute inter-rater reliability
in the PRISMA sense and does not mitigate the structural limitation described
above.
With $n=22$, the measure has insufficient statistical power to distinguish
adequate from inadequate screening consistency: the Wilson interval lower bound
of 71\% is consistent with a true intra-author agreement rate that would not
satisfy any standard IRR threshold.
The single-author screening therefore remains a limitation that readers should
weight when interpreting the corpus composition, and is discussed further in
Section~\ref{sec:discussion}.
This review used single-rater assessment throughout; independent inter-rater reliability (Cohen's $\kappa$ or Krippendorff's $\alpha$) could not be computed, and the three-pass structure, pre-specified codebook, and sensitivity analysis serve collectively as the primary procedural compensations for this limitation.
A PRISMA flow diagram summarising inclusion and exclusion counts is presented
in Figure~\ref{fig:prisma}.

% PRISMA 2020 flow diagram (TikZ)
\begin{figure}[p]
\centering
\begin{tikzpicture}[
  box/.style={
    rectangle, draw, rounded corners=3pt,
    text width=5.2cm, align=center,
    minimum height=1.0cm, font=\small
  },
  sbox/.style={
    rectangle, draw, rounded corners=3pt,
    text width=3.6cm, align=center,
    minimum height=0.9cm, font=\small, fill=gray!10
  },
  arr/.style={-{Stealth}, thick},
  lbl/.style={font=\scriptsize, text=gray}
]

% ---------- Column centres ----------
\def\cx{0}      % main column
\def\ex{6.2cm}  % exclusion column

% ---------- Row 0 : Identification header ----------
\node[box, fill=blue!10] (id) at (\cx, 0)
  {\textbf{Identification}\\[2pt]
   Database search: \textit{n} = 1\,812\\
   Snowballing: \textit{n} = 35\\
   Total: \textit{n} = 1\,847};

% ---------- Row 1 : Deduplication ----------
\node[box, below=0.7cm of id] (dedup)
  {After deduplication\\
   \textit{n} = 1\,412};

\node[sbox, right=1.4cm of dedup] (dedup_ex)
  {Duplicates removed\\
   \textit{n} = 435};

\draw[arr] (id) -- (dedup);
\draw[arr] (dedup) -- (dedup_ex);

% ---------- Row 2 : Title + abstract screening ----------
\node[box, below=0.7cm of dedup] (screen)
  {\textbf{Screening}\\[2pt]
   Title \& abstract screened\\
   \textit{n} = 1\,412};

\node[sbox, right=1.4cm of screen] (screen_ex)
  {Excluded (off-topic,\\pre-LLM era)\\
   \textit{n} = 1\,187};

\draw[arr] (dedup) -- (screen);
\draw[arr] (screen) -- (screen_ex);

% ---------- Row 3 : Full-text eligibility ----------
\node[box, below=0.7cm of screen] (elig)
  {Full-text assessed\\
   for eligibility\\
   \textit{n} = 225};

\node[sbox, right=1.4cm of elig] (elig_ex)
  {Excluded after full-text\\
   (no empirical eval,\\
   duplicate version)\\
   \textit{n} = 150};

\draw[arr] (screen) -- (elig);
\draw[arr] (elig) -- (elig_ex);

% ---------- Row 4 : Included ----------
\node[box, fill=green!10, below=0.7cm of elig] (incl)
  {\textbf{Included}\\[2pt]
   Studies in review\\
   \textit{n} = 75};

\draw[arr] (elig) -- (incl);

% ---------- Sidebar label ----------
\node[lbl, left=0.5cm of id, rotate=90, anchor=south]
  {Identification};
\node[lbl, left=0.5cm of screen, rotate=90, anchor=south]
  {Screening};
\node[lbl, left=0.5cm of elig, rotate=90, anchor=south]
  {Eligibility};
\node[lbl, left=0.5cm of incl, rotate=90, anchor=south]
  {Included};

\end{tikzpicture}
\caption{PRISMA 2020 flow diagram. Of the 1{,}847 records identified, 1{,}812
         were returned by database search (arXiv, ACL Anthology, IEEE Xplore,
         ACM DL, USENIX, NeurIPS/ICML/ICLR) and 35 were added via reference
         snowballing of included papers. Single-author screening was performed
         with adversarial re-examination of borderline cases (see
         Section~\ref{sec:methodology}).}
\label{fig:prisma}
\end{figure}


\subsection{Quality Assessment}

\paragraph{Quality appraisal instrument}
Each included primary study was appraised using a five-criterion quality instrument adapted from the Mixed Methods Appraisal Tool (MMAT) and the Joanna Briggs Institute (JBI) checklist for systematic reviews:
\begin{enumerate}
  \item Clear research question or design objective stated
  \item Study design appropriate for the research question (experimental, empirical, analytical, or formal)
  \item Evaluation methodology described with sufficient detail for replication
  \item Threat model or adversarial assumptions explicitly stated (where applicable)
  \item Limitations acknowledged by the authors
\end{enumerate}
Each criterion was scored as Met (1), Partially Met (0.5), or Not Met (0), giving a total quality score in the range [0, 5].
Studies scoring below 2.0 were flagged for heightened scrutiny during claim extraction but were not excluded from the corpus --- exclusion occurred only at the title/abstract and full-text screening phases based on topical relevance criteria, not on quality score.
Heightened scrutiny meant that claims extracted from sub-threshold studies were subjected to a mandatory additional verification note during each pass, explicitly recording whether the low quality score introduced ambiguity in the claim's reproducibility assessment; the pass/fail verdict criteria themselves were not relaxed or tightened relative to higher-scoring studies.

\paragraph{Relationship between quality ratings and the 237-claim pool}
Quality scores were not used as a threshold for claim inclusion: all quantitative and architectural assertions meeting the codebook definition (Section~\ref{subsec:verification_protocol}) were extracted from all included studies regardless of their quality score.
This choice was deliberate.
The review's primary hypothesis was that the literature's evidentiary basis is weak; excluding low-quality studies before claim extraction would have biased the verification outcome toward survival, understating the reproducibility crisis.
The five surviving claims are distributed across quality tiers: Claims~1 and~2 (safety context distillation and dual reward architecture, from Touvron et al.) derive from a study scoring 4.5/5; Claim~3 (Llama Guard architecture) from a study scoring 4.0/5; Claim~4 (Crescendo multi-turn amplification) from a study scoring 3.5/5; Claim~5 (SALAD-Bench evaluation fragility) from a study scoring 4.0/5.
The 232 refuted claims span all quality tiers, confirming that refutation rate is not primarily a function of study quality but of evaluation-protocol underspecification --- the most common failure mode (41\% of refutations).

Each included paper was also assessed on four dimensions:
(i) \textit{reproducibility}—whether experiments can be replicated from the
information provided;
(ii) \textit{generalisability}—whether results are validated on more than one model
or dataset;
(iii) \textit{adversarial rigour}—whether the paper considers adaptive attacks or
only static evaluation;
(iv) \textit{deployment relevance}—whether the proposed mechanism addresses a
realistic deployment constraint (latency, API access level, streaming).

\noindent Papers were rated Low (L) / Medium (M) / High (H) on each dimension.
Ratings were assigned after reading the full text of each paper; all four dimensions
were rated by the same single author.
As a partial calibration check, eight papers (approximately 10\% of the corpus,
selected by systematic sampling) were re-rated on all four dimensions after an
interval of four weeks without reference to the initial ratings.
Raw agreement between first and second ratings was 78\% (25/32 dimension-level
comparisons); the seven discrepancies all involved boundary cases between adjacent
categories (e.g.\ L vs.\ M), not disagreements of two categories.
This figure is reported as intra-author consistency only and cannot substitute for
inter-rater reliability.
The distribution figures in Table~\ref{tab:quality} should therefore be understood
as directional indicators of corpus-wide quality patterns, not precisely calibrated
ordinal measurements: the integers reflect the single author's consistent application
of criteria across readings, with boundary-case uncertainty of approximately
$\pm 1$--2 papers per cell.
Table~\ref{tab:quality} summarises the quality distribution across the 75
included papers.

\begin{table}[t]
  \caption{Quality assessment distribution across 75 included papers.
           Ratings: L = Low, M = Medium, H = High.
           All ratings are by a single author; see text for intra-author
           consistency check (78\% raw agreement on 10\% re-rated sample).
           Integer counts carry boundary-case uncertainty of approximately
           $\pm$1--2 papers per cell and should be treated as directional
           indicators, not precisely calibrated measurements.}
  \label{tab:quality}
  \centering
  \begin{tabular}{lrrr}
    \toprule
    \textbf{Dimension} & \textbf{L} & \textbf{M} & \textbf{H} \\
    \midrule
    Reproducibility       & 41 & 27 & 7  \\
    Generalisability      & 18 & 38 & 19 \\
    Adversarial rigour    & 33 & 29 & 13 \\
    Deployment relevance  & 22 & 34 & 19 \\
    \bottomrule
  \end{tabular}
\end{table}

The distribution confirms the reproducibility crisis: 55\% of papers receive a
Low reproducibility rating, meaning their core experiments cannot be replicated
from the information provided.
Ratings inform the confidence level assigned to specific claims in the review.

\subsection{Data Extraction}

For each included paper we extracted: proposed mechanism, attack surface addressed,
evaluation datasets and metrics, reported performance figures, latency or overhead data
(where available), and the authors' own limitations discussion.
Extraction was performed against a pre-specified codebook (eight fields per paper,
defined before any primary-source reading began) to reduce post-hoc field selection.
The codebook template is available from the corresponding author on request.

\paragraph{Corpus composition and bibliography scope}
Table~\ref{tab:corpus} enumerates all 75 primary corpus papers explicitly and
in full.
This enumeration is the authoritative corpus definition; the BibTeX file
(\texttt{main.bib}) additionally contains eleven \emph{non-corpus} references:
two SLR methodology references (Kitchenham \& Charters; PRISMA 2020); three
background/framing references (Bommasani et al.; Weidinger et al.; Perez et al.\
2022 red-teaming); four constrained-generation adjacent references (FUDGE, COLD
Decoding, NeuroLogic, speculative decoding); and two industry technical reports
cited solely to confirm the absence of reproducible audio/video safety evaluations
(Anthropic Claude~3 model card; Google Gemini technical report).
Every other entry in \texttt{main.bib} corresponds to a row in
Table~\ref{tab:corpus}.

\small
\begin{longtable}{rp{12.0cm}c}
\caption{Complete enumeration of the 75 primary corpus papers, sorted by theme.
         \textbf{Status}: P = arXiv preprint; W = workshop; C = peer-reviewed
         conference; J = peer-reviewed journal.
         C and J entries have undergone formal peer review; P entries have not.}
\label{tab:corpus} \\
\toprule
\textbf{\#} & \textbf{Reference} & \textbf{St.} \\
\midrule
\endfirsthead
\multicolumn{3}{c}{\tablename~\thetable{} (continued)} \\
\toprule
\textbf{\#} & \textbf{Reference} & \textbf{St.} \\
\midrule
\endhead
\midrule
\multicolumn{3}{r}{\itshape Continued on next page} \\
\endfoot
\bottomrule
\endlastfoot
    \multicolumn{3}{l}{\textit{Alignment foundations and intrinsic safety (4)}} \\
    1  & Bai et al.\ 2022 — Constitutional AI~\cite{bai2022constitutional} & P \\
    2  & Ouyang et al.\ 2022 — InstructGPT / RLHF~\cite{ouyang2022training} & C \\
    3  & Touvron et al.\ 2023 — Llama~2~\cite{touvron2023llama2} & P \\
    4  & Qi et al.\ 2023 — Fine-tuning safety degradation~\cite{qi2023finetuning} & P \\
    \midrule
    \multicolumn{3}{l}{\textit{Real-time and inference-time guardrails (10)}} \\
    5  & Inan et al.\ 2023 — Llama Guard~\cite{inan2023llama} & P \\
    6  & Rebedea et al.\ 2023 — NeMo Guardrails~\cite{rebedea2023nemo} & P \\
    7  & Zhang et al.\ 2024 — ShieldLM~\cite{zhang2024shieldlm} & P \\
    8  & Han et al.\ 2024 — WildGuard~\cite{han2024wildguard} & P \\
    9  & Dong et al.\ 2024 — Building Guardrails~\cite{dong2024building} & P \\
    10 & Su et al.\ 2024 — AEGIS~\cite{su2024aegis} & P \\
    11 & Ghosh et al.\ 2025 — Aegis2.0~\cite{ghosh2024aegis2} & P \\
    12 & Kumar et al.\ 2023 — Certifying LLM safety~\cite{kumar2023certifying} & P \\
    13 & Varshney et al.\ 2024 — Systematic defence evaluation~\cite{varshney2023art} & P \\
    14 & Zou et al.\ 2024 — Circuit Breakers~\cite{zou2024improving} & P \\
    \midrule
    \multicolumn{3}{l}{\textit{Jailbreak attacks (text) (20)}} \\
    15 & Zou et al.\ 2023 — GCG universal adversarial attacks~\cite{zou2023universal} & P \\
    16 & Yang et al.\ 2023 — Shadow alignment~\cite{yang2023shadow} & P \\
    17 & Wei et al.\ 2023 — Jailbroken (NeurIPS)~\cite{wei2023jailbroken} & C \\
    18 & Wei et al.\ 2023 — In-context demonstrations jailbreak~\cite{wei2023jailbreak} & P \\
    19 & Mehrotra et al.\ 2023 — Tree of Attacks~\cite{mehrotra2023tree} & P \\
    20 & Liu et al.\ 2023 — AutoDAN (stealthy)~\cite{liu2023autodan} & P \\
    21 & Zhu et al.\ 2023 — AutoDAN (interpretable)~\cite{zhu2023autodan} & P \\
    22 & Perez \& Ribeiro 2022 — Ignore previous prompt~\cite{perez2022ignore} & P \\
    23 & Shen et al.\ 2023 — ``Do Anything Now'' wild jailbreaks~\cite{shen2023anything} & P \\
    24 & Yuan et al.\ 2023 — Cipher jailbreaks~\cite{yuan2023gpt4} & P \\
    25 & Deng et al.\ 2023 — Multilingual jailbreaks~\cite{deng2023multilingual} & P \\
    26 & Ding et al.\ 2023 — Nested jailbreak prompts~\cite{ding2023wolf} & P \\
    27 & Zeng et al.\ 2024 — Persuasion-based jailbreaks~\cite{zeng2024johnny} & P \\
    28 & Mo et al.\ 2024 — Prompt adversarial tuning~\cite{mo2024fight} & P \\
    29 & Ganguli et al.\ 2022 — Red teaming~\cite{ganguli2022red} & P \\
    30 & Kim et al.\ 2024 — Shortcut reasoning jailbreaks~\cite{kim2024break} & P \\
    31 & Xhonneux et al.\ 2024 — Efficient adversarial training~\cite{xhonneux2024efficient} & P \\
    32 & Wei et al.\ 2024 — Brittleness of safety alignment~\cite{wei2024assessment} & P \\
    33 & Carlini et al.\ 2024 — Adversarial alignment~\cite{carlini2024aligned} & P \\
    34 & Wan et al.\ 2023 — Instruction-tuning poisoning~\cite{wan2023poisoning} & P \\
    \midrule
    \multicolumn{3}{l}{\textit{Defence mechanisms (text) (9)}} \\
    35 & Robey et al.\ 2023 — SmoothLLM~\cite{robey2023smoothllm} & P \\
    36 & Cao et al.\ 2023 — Robustly aligned LLM~\cite{cao2023defending} & P \\
    37 & Pisano et al.\ 2023 — Bergeron~\cite{pisano2023bergeron} & P \\
    38 & Wang et al.\ 2024 — Backtranslation defence~\cite{wang2024defending} & P \\
    39 & Xie et al.\ 2023 — Self-reminders (Nature MI)~\cite{xie2023defending} & J \\
    40 & Phute/Helbling et al.\ 2023 — LLM Self Defense~\cite{helbling2023llm} & P \\
    41 & Jain et al.\ 2023 — Baseline adversarial defences~\cite{jain2023baseline} & P \\
    42 & Goldowsky-Dill et al.\ 2023 — Localizing model behaviour~\cite{goldowsky2023localizing} & P \\
    43 & Shi et al.\ 2023 — Irrelevant context distraction~\cite{shi2023large} & P \\
    \midrule
    \multicolumn{3}{l}{\textit{Evaluation and benchmarking (8)}} \\
    44 & Li et al.\ 2024 — SALAD-Bench~\cite{li2024salad} & P \\
    45 & Mazeika et al.\ 2024 — HarmBench~\cite{mazeika2024harmbench} & P \\
    46 & Souly et al.\ 2024 — StrongREJECT~\cite{souly2024strongreject} & P \\
    47 & Zhu et al.\ 2024 — GRAM~\cite{zhu2024gram} & P \\
    48 & Zhou et al.\ 2024 — EasyJailbreak~\cite{zhou2024easyjailbreak} & P \\
    49 & Zhang et al.\ 2023 — SafetyBench~\cite{zhang2023safetybench} & P \\
    50 & Sun et al.\ 2023 — Chinese LLM safety assessment~\cite{sun2023safety} & P \\
    51 & Bhatt et al.\ 2023 — Purple Llama CyberSecEval~\cite{bhatt2023purple} & P \\
    \midrule
    \multicolumn{3}{l}{\textit{Multi-turn attacks (5)}} \\
    52 & Russinovich et al.\ 2024 — Crescendo~\cite{russinovich2024greatescalation} & P \\
    53 & Yang et al.\ 2024 — Chain of Attack~\cite{yang2024chain} & P \\
    54 & Li et al.\ 2024 — LLM robustness vs.\ jailbreak~\cite{li2024llm} & P \\
    55 & Cheng et al.\ 2024 — Multi-round interaction attacks~\cite{cheng2024leveraging} & P \\
    56 & Yi et al.\ 2023 — Indirect prompt injection~\cite{yi2023benchmarking} & P \\
    \midrule
    \multicolumn{3}{l}{\textit{Multimodal attacks and defences (17)}} \\
    57 & Gong et al.\ 2023 — FigStep typographic attacks~\cite{gong2023figstep} & P \\
    58 & Zhao et al.\ 2024 — Multimodal jailbreak evaluation~\cite{zhao2024evaluating} & P \\
    59 & Liu et al.\ 2023 — Query-relevant image attacks~\cite{liu2023queryrelevant} & P \\
    60 & Luo et al.\ 2024 — JailBreakV-28K~\cite{luo2024jailbreakv} & P \\
    61 & Niu et al.\ 2024 — Multimodal jailbreaking attack~\cite{niu2024jailbreaking} & P \\
    62 & Dong et al.\ 2024 — Patched visual prompt injection~\cite{dong2024safeguarding} & P \\
    63 & Pi et al.\ 2024 — MLLM-Protector~\cite{pi2024mllm} & P \\
    64 & Chen et al.\ 2023 — DRESS~\cite{chen2024dresssafety} & P \\
    65 & Chen et al.\ 2024 — AgentPoison~\cite{chen2024agentpoison} & P \\
    66 & Oh et al.\ 2024 — SafeVLM~\cite{oh2024safevlm} & P \\
    67 & Liu et al.\ 2024 — MM-SafetyBench~\cite{liu2024mmsafetybench} & P \\
    68 & Li et al.\ 2024 — VLGuard~\cite{li2024vlguard} & P \\
    69 & Zong et al.\ 2024 — Safety fine-tuning at no cost~\cite{zong2024safety} & P \\
    70 & Wang et al.\ 2024 — AdaShield~\cite{wang2024adashield} & P \\
    71 & Chen et al.\ 2024 — Red-teaming VLMs~\cite{chen2024red} & P \\
    72 & Bailey et al.\ 2023 — Image hijacks~\cite{bailey2023image} & P \\
    73 & Shayegani et al.\ 2023 — Jailbreak in pieces~\cite{shayegani2023jailbreak} & P \\
    \midrule
    \multicolumn{3}{l}{\textit{Industry and deployment (2)}} \\
    74 & Qi et al.\ 2024 — Visual adversarial examples~\cite{qi2024visual} & P \\
    75 & Chiang et al.\ 2024 — Chatbot Arena~\cite{chiang2024chatbot} & P \\
\end{longtable}
\normalsize

\subsection{Structured Claim Verification Protocol}
\label{subsec:verification_protocol}

The primary methodological contribution of this review is a structured
claim-verification procedure applied to all extracted claims.
The procedure uses temporally separated self-review passes by a single author to
reduce anchoring bias; it is not adversarial verification in the sense of independent
reviewers challenging each other's assessments.
The protocol is defined as follows.

\paragraph{Claim unit}
A \emph{claim} is any quantitative assertion (e.g.\ ``attack success rate reduced
from 79.8\% to 2.4\%'') or qualitative architectural assertion (e.g.\ ``the guardrail
operates on token prefixes without buffering the full response'') that the paper
presents as a result or a property of the proposed system.
Architectural claims are distinguished from quantitative claims because they require
different verification standards: architectural claims are assessed for internal
consistency and alignment with disclosed implementation details, whereas quantitative
claims are assessed for reproducibility from the provided experimental setup.

\paragraph{Verification passes}
Three temporally separated verification passes were conducted by the same author over a
period of six weeks, each separated by at least two weeks to minimise anchoring
effects.
These passes are conducted by a single author; temporal separation reduces anchoring
bias but does not constitute independent review by multiple researchers.
In each pass, the author attempted to reproduce or refute each claim by:
(i) checking whether the described experimental setup is sufficiently specified to
replicate (datasets named and versioned, models identified with version tags,
hyperparameters reported);
(ii) cross-referencing numerical results against any available code repositories
or independent reproductions in the literature;
(iii) checking for internal consistency between the claim and other figures,
tables, or proofs in the same paper.

\paragraph{Pass/fail criteria}
A pass (verdict: \textit{verified}) requires all three of the following conditions to hold in a given assessment pass:
(i) the experimental setup is sufficiently specified that a qualified reader could reproduce the reported result without contacting the authors (datasets named and versioned, models identified, hyperparameters reported);
(ii) no internal inconsistency is found between the claim and other results or proofs in the same paper;
(iii) where an independent replication or corroborating study is available in the literature, the directional finding is consistent across sources.
Condition~(i) establishes that a claim is transcribed accurately from the source and that its replication prerequisites are stated; it does not establish that the underlying experimental result has been independently reproduced or that the claim is free of benchmark contamination.
Condition~(iii) is a corroborative check only: corroboration upgrades confidence, but absence of corroboration does not, by itself, cause a claim to fail.
A claim is \textit{refuted} if any single pass finds that condition~(i) or condition~(ii) is not satisfied.
A claim is \textit{unverifiable} if the source material required to assess condition~(i) is not publicly accessible.
These criteria apply uniformly to both quantitative and architectural claims; the distinction affects how condition~(i) is operationalised (quantitative claims require numeric reproducibility; architectural claims require mechanistic reproducibility) but not the pass/fail logic itself.

\paragraph{Sensitivity analysis on the single-rater constraint}
The 98\% unverified figure rests on a single-rater audit.
To bound the uncertainty this introduces, consider the following: if borderline cases constitute at most 20\% of the corpus ($0.20 \times 237 = 47.4$, rounded to 47 claims), and if a second independent rater disagrees with the first rater on at most half of those borderline cases ($0.50 \times 47 = 23.5$, rounded to 24 claims, representing the most pessimistic plausible disagreement rate under the operationalised pass/fail criteria), then the maximum number of verdicts that could be overturned is 24.
The 20\% borderline rate is a structural assumption, not an empirically observed rate; it is stated as a worst-case bound.
The direction of bias is also bounded: the structured self-review protocol's strictness --- treating ambiguous claims as refuted rather than verified --- more likely causes under-counting of genuine findings than over-counting.
The five surviving claims should therefore be understood as a lower bound on reproducible findings under this codebook.
If all 24 were overturned from \emph{Refuted} to \emph{Verified} --- an extreme upper bound --- the verified count would rise from 5 to 29, which is 12\% of the corpus.
The central finding is therefore robust to plausible single-rater variability: even under the most pessimistic stated assumptions, the verified fraction remains below 15\%.

\paragraph{Verification basis for surviving claims}
Because condition~(iii) applies differently across the five surviving claims, we report the verification basis for each explicitly.
Claim~1 (Llama~2 dual reward model threshold): verified by condition~(i) alone --- threshold and test-set metrics are reported with sufficient precision in the source paper~\cite{touvron2023llama2}.
Claim~2 (Safety Context Distillation): verified by condition~(i) --- mechanism described at implementation level in~\cite{touvron2023llama2}.
Claim~3 (Llama Guard architecture): verified by conditions~(i) and~(iii) --- corroborated by subsequent releases and independent usage reports~\cite{meta2024llamaguard2,inan2023llama}.
Claim~4 (Crescendo multi-turn amplification): verified by conditions~(i) and~(iii) --- independent corroboration from Li~et~al.~\cite{li2024llm} using a different attack methodology.
Claim~5 (SALAD-Bench evaluation-protocol fragility): verified by conditions~(i) and~(iii) --- directionally corroborated by HarmBench~\cite{mazeika2024harmbench}, StrongREJECT~\cite{souly2024strongreject}, GRAM~\cite{zhu2024gram}, and EasyJailbreak~\cite{zhou2024easyjailbreak}.
Readers should note that Claims~1 and~2 are verified on the basis of internal reproducibility only; Claim~4 has the strongest independent corroboration of the five.

\paragraph{Voting rule}
A claim \emph{survives} verification if and only if it passes all three temporally
separated self-review passes without identified inconsistency or unreproducibility
(3-vote unanimous agreement).
A claim is \emph{refuted} if any single pass identifies a material inconsistency,
a missing prerequisite for reproduction, or a contradiction with accessible
primary sources.
Claims that cannot be assessed (e.g.\ because source material is behind a
non-public institutional repository) are marked \emph{unverifiable} and treated
equivalently to refuted for the purposes of the open-problems analysis.

\paragraph{Per-claim audit artefact}
To enable independent inspection of the verification outcomes, a structured
machine-readable artefact is submitted alongside this manuscript.
The artefact is a comma-separated spreadsheet in which each row corresponds to
one extracted claim and contains the following fields:
\texttt{claim\_id} (unique integer),
\texttt{source\_key} (BibTeX key of the primary source),
\texttt{claim\_text} (verbatim or lightly paraphrased assertion as extracted),
\texttt{claim\_type} (\textit{quantitative} or \textit{architectural}),
\texttt{pass1\_verdict} (\textit{verified} / \textit{refuted} / \textit{unverifiable}),
\texttt{pass2\_verdict},
\texttt{pass3\_verdict},
\texttt{final\_verdict} (\textit{survived} if all three passes are \textit{verified};
\textit{refuted} otherwise), and
\texttt{refutation\_rationale} (free-text explanation for any pass that returned
\textit{refuted} or \textit{unverifiable}).
The artefact covers all 237 claims extracted during the review; the five rows
with \texttt{final\_verdict = survived} correspond precisely to the five verified
findings enumerated in Section~\ref{sec:open_problems}.
The artefact is available from the corresponding author on request.
The complete codebook (eight extraction fields with operationalised criteria) and the full 237-claim extraction table (all three pass verdicts and refutation rationales per claim) are available from the corresponding author on request;
both artefacts are provided to referees and editors as part of the supplementary review package.

\paragraph{Limitation}
The use of a single author across all three passes is a limitation: although the
inter-pass delay reduces anchoring, it cannot substitute for genuinely independent
reviewers.
The per-claim artefact makes the individual assessments transparent and allows
any reader to contest specific verdicts; it does not, however, provide
independently-generated second opinions on those verdicts.

A second, distinct validity threat concerns the \emph{pre-specified codebook}
itself.
The codebook was designed before any primary-source reading began and specifies
eight extraction fields per paper.
If the codebook systematically omitted entire \emph{categories} of methodological
flaw, then no paper exhibiting those flaws would have those flaws flagged across
any of the three passes.
Three such categories are worth enumerating explicitly.
\emph{Confounded baselines}: the codebook checks whether the proposed guardrail
and the comparison baseline use compatible attack sets and judge models, but it
does not have a dedicated field capturing cases where the two systems differ on
additional silent dimensions (e.g.\ training data vintage, RLHF iteration count)
that could independently explain observed performance gaps.
\emph{Undisclosed RLHF stage details}: the codebook records whether hyperparameters
are reported but does not probe for the number of RLHF iterations, the checkpoint
selection criterion, or the composition of the human-preference dataset---all of
which can materially affect safety benchmark scores and are routinely omitted from
published papers.
\emph{Fine-tuning data contamination}: the codebook does not include a field
checking whether evaluation benchmarks used by a paper overlap with the paper's
own training or fine-tuning data, a form of contamination that inflates reported
numbers without constituting fabrication.

These three categories were not added to the codebook after the first pass, even
though related concerns emerged during reading.
The reason is methodological: adding fields post-hoc, after exposure to the data,
converts the codebook into a post-hoc instrument and introduces the same
confirmation risk the pre-specification was designed to prevent.
Extending the codebook would have required a fourth pass, and the resulting verdicts
could not be compared to the three-pass structure for earlier claims without
introducing inconsistency.
The appropriate resolution is a future independent re-extraction against a
separately designed codebook that incorporates these categories from the outset.
These three categories are flagged explicitly so that future reviewers applying the same codebook can assess whether their omission materially biases the five surviving claims;
Whether inclusion of these three categories would alter the five surviving claims cannot be determined from the current codebook alone; this is flagged as a direction for follow-on methodological work.

Temporal separation reduces anchoring bias \emph{within} the categories the
codebook captures; it provides no protection against blind spots in category
coverage.
This threat cannot be resolved post-hoc without independent re-extraction against
a separately designed codebook.
Readers should therefore treat the five verified surviving claims as the set of
claims that survived scrutiny under \emph{this codebook's extraction categories},
not as the exhaustive set of reproducible findings in the literature.
A minimum requirement for independent corroboration of the headline finding
(5~of~237 claims verified) would be for a second independent coder to apply the
same pre-specified codebook to a 20--25\% random sample of the 237 claims
($n \approx 47$--59 claims) and report pairwise agreement on both claim
extraction and pass verdicts.
Such a check would establish whether the codebook produces consistent results
across coders without requiring full re-extraction of all 237 claims.
The codebook template is available from the corresponding author to permit
external assessment of its category coverage and to enable this check.

The protocol is disclosed at this level of detail so that future SLRs in this
domain can adopt or extend it with multiple reviewers, a codebook designed by
consensus, and formally computed inter-rater agreement.
